\documentclass[12pt,a4paper]{article}
\usepackage[utf8]{inputenc}
\usepackage[T1]{fontenc}
\usepackage{amsmath,amssymb}
\usepackage{enumitem}
\usepackage{geometry}
\usepackage{array}
\usepackage{multicol}
\geometry{margin=2.2cm}
\setlist{itemsep=0.25em, topsep=0.3em}
\title{Aula 14 -- Técnicas de Recuperação de BD}
\author{Banco de Dados -- Guia de Revisão}
\date{}
\begin{document}
\maketitle
\tableofcontents
\newpage

\section{Objetivo da aula}
Esta aula trata de como o SGBD recupera o banco após falhas. O foco é manter atomicidade e durabilidade mesmo quando há queda de energia, erro de sistema, falha de transação ou crash.

\section{Ideia central}
Se uma transação confirmou, seus efeitos devem permanecer. Se abortou ou ficou incompleta, seus efeitos não devem permanecer. A recuperação usa log, checkpoints, undo e redo.

\section{Tipos de falha}
\begin{itemize}
    \item \textbf{Falha de transação}: erro lógico, violação de restrição, deadlock ou abort.
    \item \textbf{Falha do sistema}: queda do SGBD ou energia; memória principal é perdida.
    \item \textbf{Falha de mídia}: problema no disco, exigindo backup.
\end{itemize}

\section{Log}
O log registra operações importantes antes que o banco seja alterado permanentemente.

Registros típicos:
\begin{verbatim}
<START T1>
<WRITE T1, X, old, new>
<COMMIT T1>
<ABORT T1>
\end{verbatim}

\section{Write-Ahead Logging}
Regra de WAL: o registro de log deve ser gravado em armazenamento estável antes da página de dados correspondente ser gravada no banco.

Intuição: se o banco gravou uma alteração, o SGBD precisa ter informação suficiente no log para desfazer ou refazer.

\section{UNDO e REDO}
\begin{itemize}
    \item \textbf{UNDO}: desfaz efeitos de transações não confirmadas.
    \item \textbf{REDO}: refaz efeitos de transações confirmadas que talvez não tenham chegado ao disco.
\end{itemize}

\section{Exemplo visual de log}
\begin{verbatim}
<START T1>
<WRITE T1, A, 100, 80>
<START T2>
<WRITE T2, B, 50, 70>
<COMMIT T1>
-- falha do sistema
\end{verbatim}

Após a falha:
\begin{itemize}
    \item $T_1$ confirmou: pode precisar de REDO.
    \item $T_2$ não confirmou: precisa de UNDO.
\end{itemize}

\section{Checkpoint}
Checkpoint reduz o trabalho de recuperação. Ele indica um ponto em que o SGBD sincronizou informações suficientes para não precisar analisar todo o log desde o início.

Registro típico:
\begin{verbatim}
<CHECKPOINT>
\end{verbatim}

Com checkpoint, a recuperação procura transações ativas no checkpoint e transações iniciadas depois dele.

\section{Transações confirmadas e não confirmadas}
Após uma falha:
\begin{itemize}
    \item transações com \texttt{COMMIT} no log devem ser preservadas;
    \item transações sem \texttt{COMMIT} devem ser desfeitas.
\end{itemize}

\section{Deferred update e immediate update}
\subsection{Deferred update}
As atualizações só chegam ao banco depois do commit. Em geral, precisa de REDO para transações confirmadas.

\subsection{Immediate update}
As atualizações podem chegar ao banco antes do commit. Pode exigir UNDO para transações não confirmadas e REDO para confirmadas.

\section{Backup}
Falha de mídia pode exigir restauração de backup e reaplicação de log.

Estratégia geral:
\begin{enumerate}
    \item Restaurar backup.
    \item Aplicar REDO de transações confirmadas desde o backup.
    \item Garantir que transações não confirmadas não permaneçam.
\end{enumerate}

\section{Relação com ACID}
\begin{itemize}
    \item Recuperação protege \textbf{atomicidade}: transação incompleta não fica pela metade.
    \item Recuperação protege \textbf{durabilidade}: transação confirmada não desaparece.
\end{itemize}

\section{Pegadinhas comuns}
\begin{itemize}
    \item \texttt{COMMIT} no log é decisivo para saber se preserva a transação.
    \item UNDO não é para transação confirmada.
    \item REDO não é para transação abortada.
    \item Checkpoint não é backup completo; é um ponto de controle no log.
    \item Falha de mídia é mais grave que falha do sistema.
\end{itemize}

\section{Desafios}
\subsection*{Desafio 1}
Considere o log:
\begin{verbatim}
<START T1>
<WRITE T1, X, 10, 20>
<COMMIT T1>
<START T2>
<WRITE T2, Y, 5, 9>
-- falha
\end{verbatim}
O que fazer com $T_1$ e $T_2$?

\subsection*{Desafio 2}
Por que WAL é importante?

\section{Resolução dos desafios}
\textbf{Desafio 1}: $T_1$ confirmou, então deve ser preservada e pode precisar de REDO. $T_2$ não confirmou, então deve ser desfeita com UNDO.

\textbf{Desafio 2}: porque garante que o log necessário para desfazer/refazer esteja salvo antes das alterações no banco, permitindo recuperação correta.

\section{Checklist final}
\begin{itemize}
    \item Sei diferenciar UNDO e REDO?
    \item Sei analisar log com \texttt{START}, \texttt{WRITE}, \texttt{COMMIT} e falha?
    \item Sei para que serve checkpoint?
    \item Sei relacionar recuperação com atomicidade e durabilidade?
    \item Sei diferenciar falha de transação, sistema e mídia?
\end{itemize}

\end{document}
